\documentclass[12pt]{article}
\usepackage[utf8]{inputenc}
\usepackage[spanish]{babel}
\usepackage{amsmath}
\usepackage{amssymb}
\usepackage{cancel}
\usepackage{geometry}
\usepackage{booktabs}
\usepackage{array}
\geometry{margin=2.5cm}

\begin{document}
	
	\begin{center}
		{\Large \textbf{Problema 4: Perfiles Clínicos}} \\[4pt]
		{\large Operaciones de Conjuntos}
	\end{center}
	
	\bigskip
	
	\section*{Descripción del Problema}
	
	Se estudia un grupo de \textbf{20 personas} para identificar perfiles clínicos asociados
	a su bienestar emocional, cognitivo y social. El universo de estudio es:
	\[
	U = \{1, 2, 3, \ldots, 20\}
	\]
	
	A partir de entrevistas clínicas se identifican las siguientes características:
	\begin{itemize}
		\item \textbf{Ansiedad elevada} $(A)$: altos niveles de activación emocional,
		preocupación constante o síntomas ansiosos.
		\item \textbf{Dificultades en la interacción social} $(S)$: problemas en habilidades
		sociales, comunicación o relaciones interpersonales.
		\item \textbf{Estrés crónico} $(E)$: sobrecarga sostenida, fatiga emocional o
		dificultades para afrontar demandas del entorno.
	\end{itemize}
	
	\section*{Conjuntos Dados}
	
	\begin{align*}
		A &= \{2, 4, 6, 8, 10, 12, 14, 16, 18, 20\} \\[4pt]
		S &= \{5, 7, 9, 11, 13\} \\[4pt]
		E &= \{3, 6, 9, 12, 15, 18\}
	\end{align*}
	
	\noindent\textbf{Observación clave:} $A \cap S = \emptyset$, es decir, ningún elemento
	de $S$ pertenece a $A$.
	
	\bigskip
	
	\section*{Parte a) — Conjunto por Comprensión}
	
	Se determina el conjunto:
	\[
	(A - E) \cup (S \cap E)
	\]
	
	\subsection*{Paso 1: $A - E$}
	
	Se eliminan de $A$ los elementos que pertenecen a $E$
	(es decir, $6, 12, 18 \in A \cap E$):
	\[
	A - E
	= \{2,\, 4,\, \cancel{6},\, 8,\, 10,\, \cancel{12},\, 14,\, 16,\, \cancel{18},\, 20\}
	= \{2, 4, 8, 10, 14, 16, 20\}
	\]
	
	\subsection*{Paso 2: $S \cap E$}
	
	Se buscan los elementos comunes entre $S$ y $E$:
	\[
	S \cap E
	= \{5, 7, \mathbf{9}, 11, 13\} \cap \{3, 6, \mathbf{9}, 12, 15, 18\}
	= \{9\}
	\]
	
	\subsection*{Paso 3: Unión}
	
	\[
	(A - E) \cup (S \cap E)
	= \{2, 4, 8, 10, 14, 16, 20\} \cup \{9\}
	= \{2, 4, 8, 9, 10, 14, 16, 20\}
	\]
	
	\medskip
	\noindent\textbf{Descripción por comprensión:} Conjunto de personas que
	\textit{presentan ansiedad elevada pero no estrés crónico}, o que
	\textit{tienen simultáneamente dificultades sociales y estrés crónico}.
	
	\bigskip
	
	\section*{Parte b) — Conjunto por Extensión y Cardinalidad}
	
	Se determina el conjunto (aplicando precedencia estándar de operaciones):
	\[
	\bigl[(A - S)^c \cap A^c\bigr] - E
	\]
	
	\subsection*{Paso 1: $A - S$}
	
	Como $A \cap S = \emptyset$, ningún elemento de $S$ pertenece a $A$, por lo tanto:
	\[
	A - S = A = \{2, 4, 6, 8, 10, 12, 14, 16, 18, 20\}
	\]
	
	\subsection*{Paso 2: $(A - S)^c = A^c$}
	
	Se calcula el complemento respecto a $U$:
	\[
	A^c = U - A = \{1, 3, 5, 7, 9, 11, 13, 15, 17, 19\}
	\]
	
	\subsection*{Paso 3: $(A-S)^c \cap A^c$}
	
	Como $(A-S)^c = A^c$, la intersección es:
	\[
	A^c \cap A^c = A^c = \{1, 3, 5, 7, 9, 11, 13, 15, 17, 19\}
	\]
	
	\subsection*{Paso 4: $A^c - E$}
	
	Se eliminan de $A^c$ los elementos que pertenecen a $E$
	(es decir, $3, 9, 15 \in A^c \cap E$):
	\[
	A^c - E
	= \{1,\, \cancel{3},\, 5,\, 7,\, \cancel{9},\, 11,\, 13,\, \cancel{15},\, 17,\, 19\}
	= \{1, 5, 7, 11, 13, 17, 19\}
	\]
	
	\medskip
	\noindent\textbf{Resultado por extensión:}
	\[
	(A-S)^c \cap A^c - E = \{1,\, 5,\, 7,\, 11,\, 13,\, 17,\, 19\}
	\]
	
	\noindent\textbf{Cardinalidad:}
	\[
	\bigl|(A-S)^c \cap A^c - E\bigr| = 7
	\]
	
	\bigskip
	
	\section*{Resumen de Resultados}
	
	\begin{center}
		\renewcommand{\arraystretch}{1.6}
		\begin{tabular}{ccc}
			\toprule
			\textbf{Parte} & \textbf{Conjunto resultante} & \textbf{Cardinalidad} \\
			\midrule
			a) & $\{2,\, 4,\, 8,\, 9,\, 10,\, 14,\, 16,\, 20\}$ & $8$ \\
			b) & $\{1,\, 5,\, 7,\, 11,\, 13,\, 17,\, 19\}$       & $7$ \\
			\bottomrule
		\end{tabular}
	\end{center}
	
\end{document}